\section{Reverse alternating tests}\label{sec:reverse}

The undecidability argument will use exact subgroup intervals in both
orientations, and intervals of invariant subgroups in powers of a
nonabelian simple group.  We first prove the reverse exclusion needed for
that reduction.  Its hypothesis is an arbitrary exact realization of a
reversed alternating subgroup lattice.  From such a realization we recover
an ordinary site interval, build the linear charts of the path argument
from actual elementary subgroups, and reproduce the structural exclusions.
The forward marker theorem in the next section has a complementary
conclusion: after a bounded prefix, it supplies an actual alternating
simple subgroup with many point-stabilizer labels.

In \Cref{sec:weighted}, markers will be placed beside units carrying grids
and full permutation groups.  Normal and projection tests on nontrivial
lattice images of that template will allow these two theorems to force a
forward product configuration.  Intersections of cuts then identify one
common partition into actual orbit cells, in which equal totals and
squares supply arithmetic.  Finally, \Cref{sec:counting} bounds the number
of lattice elements computably in the number of units, independently of
the formal alphabet sizes.  A computable bound for carriers of realizable
lattices would consequently bound positive Diophantine solutions whenever
they exist.

Throughout, reversal means reversal of the entire displayed parameter
lattice.  Selecting its coordinate-product labels produces a separate
ordinary site interval; it need not preserve all labels of the original
invariant interval.  We continue to use property~\((*)\): every normal
subgroup of the upper endpoint is contained in, or supplemented by, the
base, and the quotient by the base core has a unique minimal nontrivial
normal subgroup, which is nonabelian.  Fixed absolute thresholds need not
be effectively evaluated: they may be fixed constants in an implication
between the existence of algorithms.  All groups in this section are
finite.

\begin{theorem}[Reverse alternating exclusion]\label{r:reverse-exclusion}
For all sufficiently large integers $v$, the lattice
$\operatorname{Sub}(A_v)^{\mathrm{op}}$ has no exact realization as
\[
 [P,S^b]^\Delta,
\]
where $S$ is nonabelian simple, $\Delta$ acts by automorphisms transitively
on the $b$ simple factors, $P$ is a $\Delta$-invariant coordinate product,
and either $P\ne1$ or $b=1$.
\end{theorem}

\subsection{Normal tests and the ordinary site interval}

We first establish the normal test needed at every prepared link.  The
test applies to the full invariant interval through a semidirect product;
we will then transfer its conclusion to the ordinary interval formed by
the coordinate-product labels.

\begin{lemma}[Semiregular alternating normal test]\label{r:alternating-normal}
Let $X\le A_n$ be semiregular and elementary abelian, allowing $X=1$.
For $n$ sufficiently large in terms of $|X|$, every exact finite group
interval isomorphic to $[X,A_n]$ or $[X,A_n]^{\mathrm{op}}$ has
property~\((*)\).
\end{lemma}

\begin{proof}
We first show that a normal subgroup can give only an endpoint label.
The unique nonabelian monolith will then follow from nonmodularity.
Suppose $[H_0,G_0]$ is such an interval, and let $D\in[X,A_n]$ label
$H_0N$, where $N\trianglelefteq G_0$.  This label satisfies the
left-modular identity, in either orientation.  Moreover, projecting
$H_0N$ to $[L,K]\subseteq[H_0,G_0]$, by meeting and joining in either
order, gives
\[
 L(N\cap K),\qquad N\cap K\trianglelefteq K.
\]
The identity
\begin{equation}\label{r:left-modular}
 z\vee(d\wedge y)=(z\vee d)\wedge y\qquad(z\le y)
\end{equation}
is self-dual.  Consequently neither $D$ nor any of these projections
can admit a test
\begin{equation}\label{r:forbidden-test}
 D\cap y\le z<y\le\langle D,z\rangle.
\end{equation}
We show that $D=X$ or $D=A_n$.

Choose $37$ points in distinct $X$-orbits.  Their $X$-translated
$37$-tuples are pairwise disjoint.  Let $T$ be the product of the
alternating groups supported on these tuples.  Thus $X\cap T=1$, and
$[X,XT]$ identifies with the $X$-invariant subgroups of $T$.  For
$L\le A_{37}$ let $T_L$ denote its product in the matched coordinates.
Put $D_0=D\cap T$, and write $V\le A_{37}$ for a coordinate projection
of $D_0$.  The following pair of conditions is impossible whenever
$B_0\le a<c\le A_{37}$:
\begin{equation}\label{r:local-test}
 \langle V,B_0\rangle\cap c\le a,
 \qquad
 c\le\langle a^{\langle V,a\rangle}\rangle.
\end{equation}
Indeed project once more to
$D'=\langle D_0,T_{B_0}\rangle$.  The first condition gives
$D'\cap T_c\le T_a$.  Conjugating independently supported copies of
$a$ by elements whose site projection belongs to $V$ gives
$T_c\le\langle D',T_a\rangle$, contrary to
\eqref{r:forbidden-test}.

Suppose first that $V$ is proper and transitive.  Burnside's prime-degree
Theorem \citep{Muller2005} and the classification of finite
two-transitive groups imply
that $V$ lies in the affine normalizer of a regular cyclic group of
order $37$; in particular it contains a $37$-cycle.  Here the
non-affine possibilities at degree $37$ all contain $A_{37}$.
For completeness, $37$ is not a projective point number
$1+s+\cdots+s^j$ for a prime power $s$: for $j=1$ one would have
$s=36$, for $j=2$ the adjacent possibilities $s=5,6$ give $31,43$,
and for $j\ge3$ the small remaining possibilities likewise miss $37$.
The other classical and exceptional two-transitive degrees give none
at $37$; we use the permutation-group classification in
\citet[Table~1]{MeagherSpigaTiep2016}.

Choose a cyclic $37$-group $a$ so that
$c=N_{A_{37}}(a)$ intersects $V$ trivially.  To justify the choice, both
$V$ and $c$ have order at most $37\cdot18=666$.  A nonidentity affine
permutation fixes at most one point.  From its cycles one can choose
nine disjoint directed moved pairs.  A commuting permutation is
determined on those eighteen points by the ordered images of the nine
sources, and has at most $19!$ choices on the rest.  Its centralizer
in $S_{37}$ therefore has order at most
$(37\cdot36\cdots29)19!$, so its $A_{37}$-class has size at least
$28!/(2\cdot19!)>666^2$.  Averaging the intersections of $V$ with
conjugates of $c$ now gives a conjugate with no nonidentity intersection.
The groups $V$ and $a$ contain distinct regular cyclic $37$-groups;
a proper transitive subgroup would be affine, where the Sylow
$37$-group is unique.  Hence $\langle V,a\rangle=A_{37}$.
Its simplicity gives \eqref{r:local-test} with $B_0=1$, a contradiction.

Suppose next that $V$ is nontrivial and intransitive.  Choose
complementary invariant sets $A,F$ as follows: if there are at least two
singleton orbits, use two of them as $F$; if there is exactly one, use
that one; if there are none, use a smallest orbit.  In all cases
$|A|\ge19$.  Remove two points from $A$ when $|F|\ge2$, and three
when $|F|=1$, leaving a set $Z$ meeting every orbit in $A$.  In the
first case with two fixed points removed to form $F$, this is possible
because a nontrivial even action has
$\sum_O(|O|-1)\ge2$ on its other orbits.  In the remaining cases the
orbits of $A$ are nonsingletons, so the required removals are possible.
Set $B_0=A_Z$.  The $V$-conjugates of $B_0$ generate $A_A$:
their supports cover $A$ and have large overlaps.

Take a $4$-cycle on the removed points and enough points of $F$ to
make four.  When two points come from each set, arrange the two from
$A$ consecutively and the two from $F$ consecutively.  Multiply this
cycle by a disjoint transposition in $Z$, obtaining an even permutation
$g$.  Set
\[
 a=\langle B_0,g^2\rangle,\qquad c=\langle B_0,g\rangle.
\]
No element of $c\setminus a$ preserves $A$, whereas
$\langle V,B_0\rangle$ does.  Thus the first condition in
\eqref{r:local-test} holds.  The normal closure of $a$ in
$\langle V,a\rangle$ contains $A_A$ and its conjugate by $g^2$.
Their overlapping supports generate the alternating group on $A$
with the used points of $F$ added.  This group contains $c$, proving
the second condition and yielding another contradiction.

The two exclusions show that every site projection $V$ is $1$ or
$A_{37}$.  We next rule out diagonal ties in the latter case.
Then $D_0$ is subdirect and in fact equals $T$.
Otherwise the description of subdirect simple powers as products of
twisted diagonals gives a tied pair of sites.  Choose
$a=C_2<c=C_6\le A_{37}$.  Independent nontrivial supports from
$T_a$ break all ties, so $\langle D_0,T_a\rangle=T$.
Modularity would give
\[
 T_c=\langle D_0\cap T_c,T_a\rangle.
\]
On a tied pair, quotient each $C_6$ coordinate by its $C_2$ subgroup.
The image of $D_0\cap T_c$ has size at most $6$, whereas that of the
left side has size $3^2=9$.  The displayed equality is impossible.
Thus every transversal product $T$ is either contained in $D$ or has
trivial intersection with $D$.

If one such $T$ is contained in $D$, then all are.  Change one point
of a transversal at a time, using spare $X$-orbits when needed.
Consecutive transversals share supported three-cycles, so the
alternative of trivial intersection is impossible at the next one.
The supported $A_{37}$'s on the transversals themselves have connected,
overlapping supports covering the letter set.  They generate $A_n$,
and hence $D=A_n$.

Otherwise $D\cap T=1$ for every transversal product, also for shorter
transversals after extending them to length $37$.  We will show that every
element of $D$ then belongs to $X$.  First suppose $d\in D$
sends a point outside its $X$-orbit.  Consider the directed non-loop
moves on the set of $X$-orbits.  If three have six distinct endpoint
orbits, use a three-cycle $u$ on their sources; its $d$-conjugate
$u'$ is supported on their destinations.  If not, endpoints of a
maximal disjoint matching comprise at most four designated orbits,
and every non-loop move meets one of them.  Since each designated
orbit has size $|X|$, only boundedly many orbits, in terms of $|X|$,
are touched.  On sufficiently many untouched orbits, two points in
different orbits have the same $X$-offset under $d$.  Multiplying $d$
by an element of $X\le D$, make it fix those two points.  Use a
three-cycle $u$ on them and one point moved outside its orbit, and
again put $u'=u^d$.  In either case one transversal product $T$
contains $u,u'$, and
\[
 z=XT_{\langle u\rangle}<y=\langle z,u'\rangle\le XT,
 \qquad D\cap y=X,\qquad y\le\langle D,z\rangle.
\]
This contradicts \eqref{r:forbidden-test}.

It remains to exclude $d\in D\setminus X$ preserving each $X$-orbit.
Some transversal five-tuple has its matching of $X$-translated tuples
disrupted by $d$; preservation of every such matching would force all
offsets to be one common element of $X$.  On this matching use
\[
 z=XT_{C_5}<y=XT_{A_5}.
\]
Inside $\langle D,z\rangle$, take the supported five-cycles on the
original matching, their $d$-conjugates, and all $X$-translates.
The connected components of their support hypergraph on the five
orbits have size at least ten: the original matching partitions
these orbits into five-tuples, a disrupted tuple joins at least two,
and the $X$-action is transitive on the original matching tuples.
On each component the generated group is primitive.  Indeed a
supported prime cycle cannot permute nontrivial blocks: if it moved
blocks of size at least two, it would move at least ten points.  Its
support must therefore be within a single block, and connectedness
of the support hypergraph leaves no nontrivial block system.
Jordan's prime-cycle Theorem \citep[Theorem~1.1]{Jones2014} now supplies the full alternating group
on each component, independently, because the generating cycles
are independently supported.  It follows that
$y\le\langle D,z\rangle$, again contradicting
\eqref{r:forbidden-test}.  Thus $D=X$ or $A_n$, proving the normal
endpoint assertion.

Pass to the quotient by $\operatorname{core}_{G_0}(H_0)$.
Every nontrivial minimal normal subgroup supplements the base.
An abelian one would identify the interval, by intersection, with
an invariant-subgroup lattice of an elementary abelian group, hence
a modular lattice.  But $[X,A_n]$ is nonmodular, as already witnessed
inside the transversal product labels by the intransitive test above.
If two distinct minimal normals occurred, they would commute.  Using
one as the supplement identifies the interval with its invariant
subgroups; supplementation by the other means that the base induces
all inner actions on the first.  All these intersections would then
be normal subproducts of its simple factors, again giving a modular
lattice.  The same contradictions hold for the reversed lattice.
The core quotient therefore has a unique nonabelian minimal normal
subgroup.
\end{proof}

\begin{lemma}[Product labels and the site normal test]\label{r:site-normal}
Suppose the realization in \Cref{r:reverse-exclusion} exists.
After restricting its parameter lattice to a sufficiently large
supported $A_n\le A_v$, coordinate-product labels form an ordinary
interval
\[
 [H,M],\qquad H=\Lambda Q,\quad M=\Lambda S,
\]
where $Q$ is the base projection at a fixed site and
$\Lambda\le\operatorname{Aut}(S)$ is the site-stabilizer action.
They are closed under meets and joins in the original interval.
Write $K_X$ for the ordinary label indexed by a prime-generated
$X\le A_n$.  For every bounded-size semiregular elementary $q$-group
$X$, including $1$, with $q>5$ prime, the interval $[H,K_X]$ has
property~\((*)\), has base equal to a meet of strictly intermediate
coatoms, and has arbitrarily large residual height when $n/|X|$
is sufficiently large.  These conclusions persist under the core
and centralizer quotients whose kernels lie in the base.
\end{lemma}

\begin{proof}
A proper subdirect label cannot contain a nontrivial invariant product
$P$: factor transitivity makes every base projection nontrivial, and
its independent supports normally generate each full simple site.
For $b=1$ the assertion is automatic.  A coatom therefore has proper
product closure and, by maximality, equals that closure.  A subgroup
of $A_v$ generated by prime-order elements is a join of atoms in
parameter order; its reversed label is a meet of coatoms, hence a
product.  In particular a supported $A_n$ is a product label.  It
can be used as the new base; we keep the name $P$.

Identify $S$ with its inner automorphisms and set $I=\Lambda\cap S$.
A proper coatom product projects to a nontrivial proper group $T$
normalized by $I$, maximal in the local proper invariant interval.
Since $TI=S$ would make $T\trianglelefteq S$, we have $TI<S$, whence
$TI=T$.  Intersecting these coatoms gives $I\le Q$, since the
corresponding prime-order groups generate $A_n$.  Consequently an
overgroup of $H$ in $M$ is determined by its intersection with $S$;
that intersection transports equivariantly to every site.  Conversely
each invariant site projection containing $Q$ gives such an overgroup.
Coordinate products are closed under intersection and generation,
proving the interval assertion.

Fix $X$ as in the statement and put $Z=K_X\cap S$.
The exact invariant interval from $P$ to its product label is an
ordinary interval with endpoints
\begin{equation}\label{r:semidirect-test}
 P\rtimes\Delta,\qquad (\textstyle\prod Z)\rtimes\Delta,
\end{equation}
in matched-site notation; replace $\Delta$ by its finite acting image
if necessary.  It has reversed parameter lattice $[X,A_n]$.
A $K_X$-normal subgroup $N\le Z$ transports well-definedly, by
site-stabilizer invariance, to a normal product at the upper endpoint
of \eqref{r:semidirect-test}.  By \Cref{r:alternating-normal}, unless
$N\le Q$, it supplements the lower endpoint.  Thus $QN=Z$, and
$HN=K_X$.

We also need actual product coatoms with meet $H$.  On three points
in distinct free $X$-orbits, the product of translated supported
$C_3$'s is elementary abelian and disjoint from $X$.  Coprimality
$q\ne3$ splits it into irreducible $X$-modules.  Adjoining one
nonzero irreducible summand to $X$ covers $X$, remains strictly
below $A_n$, and is prime-generated.  The joins of all these groups
contain every supported transversal three-cycle, hence generate
$A_n$.  On reversing, their product labels are strictly intermediate
coatoms whose meet is $H$.  Using disjoint triples of orbits gives
arbitrarily long chains of such prime-generated product labels.

Now let $R\trianglelefteq K_X$ satisfy $R\cap Z\le H$.
Then $[R,Z]\le H\cap Z$.  Therefore $R$ normalizes the intersection
with $Z$ of each indicated coatom.  If $R$ is not contained in one
of them, it and that coatom generate $K_X$, so the coatom's inner
intersection is normal in $K_X$.  That intersection properly contains
$H\cap Z$, but its product with $H$ is precisely the proper coatom.
This contradicts the preceding normal test for normals within $Z$.
Hence $R\le H$.

Every missing normal thus has a missing intersection with $Z$.
Two missing normals commuting modulo the base core, including a
single one abelian modulo the core, would transport their inner
intersections to missing normal products in
\eqref{r:semidirect-test}.  They still commute modulo the transported
product of the core's intersection with $Z$.  That product is normal
at the upper endpoint and lies in its base.  This contradicts the
unique nonabelian monolith from \Cref{r:alternating-normal}.
All assertions of property~\((*)\) follow.  Factoring a core or
centralizer kernel already in the base preserves the exact intervals,
the coatom meets, and these conclusions.
\end{proof}

\subsection{Elementary charts and supported normal products}

The site interval need not retain every label of the original invariant
interval.  We therefore build charts for which all comparable labels are
products, and identify the precise replacements for the cross-chart and
normal-product arguments used in the earlier pruning construction.

\begin{lemma}[Elementary charts with projective separation]
\label{r:projective-charts}
Fix a prime $q>5$, a bounded path depth $D$, a maximum increment $r$,
and finite chart dimensions and numbers.  There are arbitrarily large
integers $n$ for which one can assign semiregular elementary
$q$-subgroups of $A_n$ to all
nodes and private charts in the finite recursive construction,
compatibly at their attachments.  At a node $W$, each private chart is
$W\oplus\mathbf F_q^N$ and extends the prescribed action of $W$.
Two nonidentity root-coordinate translations from distinct private
charts attached at that same $W$ generate a common
$\operatorname{PSL}_2(q^u)$ factor.
Here $u$ may be chosen a sufficiently large odd prime at each
attachment.
\end{lemma}

\begin{proof}
At a stage over $W$ take a replicated regular action of
\begin{equation}\label{r:attachment-group}
 W\times\operatorname{PSL}_2(\mathbf F_{q^u}).
\end{equation}
Assign different projective points $e$ to the private charts, and
embed their new coordinates as $N$-dimensional $\mathbf F_q$-subspaces
in the corresponding unipotent root groups.  For distinct $e,f$,
unipotent determinant-one lifts with coordinates $x,y$ have product
trace
\[
 2+c_{ef}xy,\qquad c_{ef}\ne0.
\]
Choose the embeddings so that, for every pair of nonzero chart
vectors from different charts, this trace has square outside
$\mathbf F_q$.  Such a choice exists: independent random linear
maps from the fixed-dimensional chart spaces to $\mathbf F_{q^u}$
are injective with probability tending to one.  Nonzero input vectors
in different charts have independent uniform images.  For each
fixed nonzero $c_{ef}$ the probability that
$2+c_{ef}xy\in\mathbf F_q$ tends to zero, uniformly in the chosen
points.  There are only finitely many tests.  Since $u$ is odd, an
element of $\mathbf F_{q^u}$ whose square belongs to $\mathbf F_q$
already belongs to $\mathbf F_q$; thus the stated squared-trace
avoidance follows.

The Dickson subgroup Theorem for $\operatorname{PSL}_2$
\citep[Theorem~2.2]{Giudici2007} now implies
that the two translations generate the full projective group.
Their distinct fixed points exclude the point-stabilizer alternative;
their order $q$ excludes the dihedral groups, whose orders are
$q^u-1$ and $q^u+1$, and $q>5$ excludes the small groups
$A_4,S_4,A_5$.  Since $u$ is odd, the square-field
$\operatorname{PGL}_2$ alternative cannot occur.  Since both $q$ and
$u$ are prime, the only remaining proper subfield group is
$\operatorname{PSL}_2(q)$.  Squared traces of determinant-one lifts
in a conjugate of that group belong to $\mathbf F_q$, contrary to our choice.
The square is unchanged by changing the lift projectively.

Choose $n$ divisible by all orders needed for this finite recursive
construction and sufficiently large for every residual-height test.
At a new attachment, the given semiregular $W$-action and the
restriction of the replicated regular action
\eqref{r:attachment-group} are both disjoint unions of regular
$W$-sets of the same total size.  An equivariant bijection identifies
them, extending the existing action.  The actions lie in $A_n$:
$W$ has odd order, and the projective factor is perfect, so neither
has a nontrivial sign character.  No disjointness between later
subtrees is needed.  The recursive construction is finite, and the
orders and sizes can all be accommodated by the final choice of $n$.
\end{proof}

\begin{lemma}[Linear failure patterns]\label{r:linear-patterns}
Consider a pruning round at an elementary parameter $W$ in the charts
of \Cref{r:projective-charts}, with actual ancestor socles as in
\Cref{c:ancestor-socles}.  Assume $K_W$ avoids the final exclusion list
relative to every strict ancestor.  In a failed chart every trial
extension $U$ of the prescribed increment has a chosen proper
full-stabilizer witness at some tested ancestor $W'$, with parameter
\[
 W'\le X(U)\le U,\qquad X(U)\not\le W.
\]
In an anchored round, the available elementary extension $B$ has
$K_B$ previous-good and $\dim(B/W)$ greater than the trial increment.
Take $q$ above the absolute line bounds for the initial exclusions,
the available dimensions sufficiently large for the prescribed finite
patterns, and $n/|X|$ sufficiently large for \Cref{r:site-normal} at
every elementary parameter used.

Then the graph construction of \Cref{c:failed-pattern} supplies lines
and bounded independent patterns in every failed chart.  For two matched patterns
in distinct private charts attached at the same current $W$, there is
one ordinary subgroup $J$ such that every cross-pair intersection is
$J$, with $H\le J$, and $J$ lies in every witness label of both patterns.
Every comparable chart interval is an exact reversed linear-subspace
interval, and every tested witness and previous-good anchor has
property~\((*)\).  In the respective high-type classical and alternating
configurations, this replacement for a cross-chart intersection equal to
$H$ makes the initial exclusions apply using only a bounded line or
independent list.
\end{lemma}

\begin{proof}
All intermediate subgroups between comparable elementary chart
parameters are elementary and prime-generated.  By
\Cref{r:site-normal}, their labels are products, so these intervals
are exact even after passage to the ordinary site interval.  The
same is true between an ancestor and a comparable trial extension,
and after the indicated ancestor core or centralizer quotients.
Thus full stabilizers witnessing failure have labels $K_{X_i}$
with all prepared-link normal tests available.

Within an available extension of $W$, choose horizontal and vertical
complements from the new root coordinates.  Color graphs of linear
maps by the failure's ancestor, category, bounded case data, and
horizontal witness subspace.  The multidimensional Hales--Jewett Theorem supplies arbitrarily many affine parameter slots with all
this data fixed \citep{HalesJewett1963}.  The fixed horizontal
witness is not contained in $W$.  Choose a horizontal functional
nonzero on its image outside $W$, and vary each slot by multiples
of that functional.  Distinct parameters give distinct witness
subspaces $X_i$.  The span of two includes all witnesses on their
parameter line and adds exactly one dimension; an affinely
independent list adds a dimension at every step.  Upon reversal,
pair intersections lie below every label on that line, are covered
by the two labels, and bounded independent lists give strictly
decreasing successive intersections.  In particular all pairs of
distinct labels on one parameter line have the same intersection $E$,
and $E$ lies in every label on that line.

For two private charts attached at $W$, write $L_0,L_1\le W$ for
the projections of their witnesses to $W$.  These projections are
constant within their respective patterns, because only the vertical
root coordinates vary.  Put $G=\operatorname{PSL}_2(q^u)$ for the
attachment factor.  Any cross pair generates a subgroup of
$(L_0+L_1)\times G$ projecting onto $G$, by
\Cref{r:projective-charts}.  Its commutator subgroup lies in $G$ and
projects onto $[G,G]=G$, since $W$ is abelian.  Hence the generated
subgroup is exactly $(L_0+L_1)\times G$.  It is independent of the
chosen pair, contains every witness parameter in both patterns, and
is prime-generated.  Its ordinary reversed label $J$ is therefore
the common cross-pair intersection, and
\[
 H\le J\le K_{X_i}
\]
for every witness in either pattern.

A witness need only contain the subspace of its tested ancestor
and extend outside the current $W$; it need not contain $W$.
The fixed $W$-projections suffice for the preceding argument.
We now verify the uses of $J$ in the initial exclusions.  For subset
stabilizers, $J$ contains the supported alternating groups on the at
most four intersection cells, so it has boundedly many point orbits;
every subset in either pattern is $J$-invariant.  For missing
equal-block partitions, supplementation still gives $H$ its rich
actions on both block lists.  Thus $H$ has an absolute bound on
orbits of ordered pairs of intersection cells.  The supported
alternating groups within those cells lie in $J$; a cell-pair
stabilizer therefore has an absolute bound on point-pair orbits,
including coincident cells and cells of size at most two.  Hence $J$
has boundedly many point-pair orbits and boundedly many invariant
equivalences.  This gives the contradiction in
\Cref{c:alternating-primitivity}; the retained-block cases use the
original normal tests.  In the classical small-space cases, the
supported large group on the common complement lies in $J$ and
preserves every structure in both patterns.  The bounded-span and
linear-constraint arguments of
\Cref{c:linear-subspaces,c:isometry-subspaces} therefore apply to a
bounded independent list.

The remaining initial classical cases need only a parameter line.
For a descending list of distinct subspaces or signed flags, the
intersection of the first two stabilizers fails to preserve the
third, by the elementary maps in \Cref{c:linear-subspaces}.  This
contradicts $E\le K_{X_3}$.  For singular spaces, first discard the
at most two exceptional incomparable maximal spaces, then choose
three members of the remaining chain.  For mutually perpendicular
retained large spaces or orthogonal idempotents, choose four distinct
members.  After the first two structures are imposed, the third
cannot exhaust the unsplit remainder because the fourth nonzero
space remains.  The supported classical or special linear group on
that remainder gives the same strict third decrease.  Thus a fixed
number of labels on a line replaces every height-sized initial
comparison.
\end{proof}

We next recover the independently supported groups needed when a
normal simple power is missing.  We cannot assume that every
intermediate intersection with that power is a coordinate product:
the ordinary site interval may omit labels of the full invariant
interval.  The structural exclusions need only a previous-good
common intersection containing enough independently supported normal
subgroups.  We find it by bounding the codimension of product closures
and then selecting compatible affine parameters.  The selection uses
the following finite-vector-space observation.

\begin{lemma}[Simultaneous affine kernels]\label{r:affine-kernels}
Fix a finite field $\mathbf F_q$ and integers $h,m\ge0$.
For sufficiently large $d$, any family of linear maps
$\phi_i:\mathbf F_q^d\longrightarrow\mathbf F_q^h$, indexed by
$i\in\mathbf F_q^d$, has an affine $m$-space $A$ such that
\[
 \phi_i(j-i)=0\qquad(i,j\in A).
\]
\end{lemma}

\begin{proof}
Color every two-dimensional subspace, with its canonical reduced
row-echelon basis $(b_1,b_2)$, by the finite table
\[
 \bigl(\phi_{a b_1+b b_2}(c b_1+d b_2)\bigr)_{a,b,c,d\in\mathbf F_q}.
\]
The number of colors depends only on $q,h$.  The finite-vector-space Ramsey Theorem gives a homogeneous subspace of dimension $m+2$
\citep{GrahamLeebRothschild1972,GrahamLeebRothschild1973}.
Let $e_1,e_2,\ldots,e_{m+2}$ be its canonical basis, and put
\[
 A=e_1+\langle e_3,\ldots,e_{m+2}\rangle.
\]
Reduced row-echelon coordinates are compatible with this passage to
a subspace: multiplying a reduced row-echelon matrix in these
coordinates by the reduced row-echelon matrix of the containing
space again gives reduced row-echelon rows in ambient coordinates.
For $i\in A$ and $w\in\langle e_3,\ldots,e_{m+2}\rangle$, both
$(i,e_2)$ and $(i,e_2+w)$ are such canonical plane bases.
Their colors agree, so $\phi_i(e_2)=\phi_i(e_2+w)$.
Linearity gives $\phi_i(w)=0$, which proves the claim.
\end{proof}

\begin{lemma}[Supported products in a selected pattern]
\label{r:supported-products}
Let $X_i\le B$ be the elementary parameters of a failure pattern in
\Cref{r:linear-patterns}, with $K_B$ a previous-good common anchor.
Suppose the witness
$K_{X_i}$ has a missing normal simple power $Y_i$, possibly modulo
a retained normal subgroup, and $H$ is transitive on its simple sites
with an absolute bound on its ordered-pair orbit count.  After
selecting a line, or an affine subpattern of any prescribed positive
dimension, its common intersection label $E$ has property~\((*)\),
contains $K_B$ and hence remains previous-good, and, for each $i$,
contains an $E$-normal product
\[
 N_i=\prod_a R_{i,a}\ \le E\cap Y_i
\]
of independently supported proper subgroups, not contained in $H$.
Its image in $E/\operatorname{core}_E(H)$ contains the monolith.
The supported factors are nonsoluble, and the same statement with
preimages applies when retained centers have been factored out.
\end{lemma}

\begin{proof}
Write $Y=Y_i$ temporarily, so $HY=K_{X_i}$.
Every intermediate overgroup is $H$ times its intersection with $Y$.
Take the coordinate-product closure of this intersection and multiply
by $H$.  For $X_i\le Z\le B$ this is an overgroup
$K_{g_i(Z)}$, where
\[
 X_i\le g_i(Z)\le Z,
\]
and $g_i$ is monotone.  The resulting group's intersection with $Y$
is the product closure itself.  Therefore $g_i(Z)=X_i$ forces the
intersection at $K_Z$ to be subdirect in $Y$.  All intermediate
intersections from $K_Z$ to $K_{X_i}$ are then subdirect.  Their
tie relations are unions of $H$-orbits on ordered pairs of sites;
a strict inclusion forces strict refinement of ties, because
comparable subdirects with the same ties coincide.  The pair-orbit
bound gives an absolute $h$ with
\[
 g_i(Z)=X_i\quad\Longrightarrow\quad\dim(Z/X_i)\le h.
\]
Choose a complement $C$ of $g_i(B)/X_i$ in $B/X_i$.
Monotonicity gives $g_i(C)\le g_i(B)\cap C=X_i$ (using the preimage
of $C$ in $B$ in this notation), so
\begin{equation}\label{r:closure-codimension}
 \operatorname{codim}_B g_i(B)\le h.
\end{equation}
Exactness between comparable elementary labels ensures that every
closure label used here really is indexed by the indicated subspace.
The same calculation can be made modulo a retained normal subgroup;
taking its preimages does not alter the intervals.

In the graph pattern, $X_j$ differs from $X_i$ in the single vertical
direction determined linearly by the parameter difference $j-i$.
The preimage of $g_i(B)$ in these vertical directions has codimension
at most $h$, by \eqref{r:closure-codimension}.  Choose a map $\phi_i$
to $\mathbf F_q^h$ with that preimage as kernel.
Apply \Cref{r:affine-kernels}, asking beforehand for enough
Hales--Jewett slots.  In the resulting affine subpattern,
$X_j\le g_i(B)$ for every $i,j$.  Let $U$ be the span of its witnesses
and put $E=K_U$.  Then $U\le g_i(B)$ for all $i$, and consequently
\[
 H<K_{g_i(B)}\le E<K_{X_i}.
\]
Here $K_{g_i(B)}>H$ because $g_i(B)\le B<A_n$.
The label $E$ is above $K_B$, so it inherits all previous-good
exclusions, and \Cref{r:site-normal} gives property~\((*)\) for
$[H,E]$.

The intersection $K_{g_i(B)}\cap Y_i$ is a coordinate product not
contained in the base; otherwise supplementation would give
$K_{g_i(B)}=H$.  Each of its single-site components already lies
independently in $E\cap Y_i$.  Let $R_{i,a}$ be the full subgroup
of $E\cap Y_i$ supported on site $a$ and put
$N_i=\prod_aR_{i,a}$.  Conjugation by $E$ permutes these support
subgroups, so $N_i\trianglelefteq E$, and $N_i\not\le H$.
Site transitivity and $E<K_{X_i}$ imply that each factor is proper:
one full simple factor would force all of them and then $E=K_{X_i}$.
The factors are mutually conjugate under $H$, and if they were
soluble, $N_i$ would be a soluble missing normal, contrary to
property~\((*)\).  Thus they are nonsoluble.  The image of $N_i$
contains the common nonabelian monolith of
$E/\operatorname{core}_E(H)$, again by property~\((*)\).

We spell out the scope needed in the structure exclusions.
For direct systems calculated modulo block centers, each block
carries at most two simple sites, and the at-most-two block classes
have independent alternating top actions induced already by $H$
through supplementation; in linear type there is one list.
These give the required absolute pair-orbit bound.  In the case
of several site orbits one instead uses the supports of a retained
orbit as in \Cref{c:supported-anchor-groups}.  Lifting independent factors on a
block to its natural supported matrix group gives projective image
in $E$, since the lost centers were retained.  The lift is preserved
by the semilinear block stabilizer in $E^+$ and supplies the groups
of order greater than two required in the block-module argument.
If needed take products of the supported factors over one block.

For tensor and Cartesian systems the $Y_i$ already have independent
simple supports.  The common-monolith and iterated-commutator
arguments of \Cref{c:tensor-refinement,c:cartesian-refinement} apply to these
$E$-normal products $N_i$, without requiring
$N_i=E\cap Y_i$.  Their tuple intersections still generate a
nontrivial $E$-normalized subgroup.  It spans in the tensor case,
and is transitive in the Cartesian case, by the previous-good
exclusions inherited from $K_B$.  In the diagonal primitive case
choose a minimal $E$-normal subgroup inside the nontrivial product
of proper supported factors.  This is exactly the containment
needed by that argument.  Retained centers, when present, are
normal in $E$ and lie in $H$, so these conclusions survive their
removal and reinstatement.
\end{proof}

\begin{proof}[Proof of \Cref{r:reverse-exclusion}]
Prepare the charts in \Cref{r:projective-charts}, and use
\Cref{r:site-normal} for every elementary parameter encountered.
At every node the current almost simple quotient is the one from
\Cref{c:ancestor-socles}; initially it is $M$.  The normal, residual-height,
and exact comparable-link hypotheses have all been verified.
\Cref{r:linear-patterns} supplies the line, independent-list, and
cross-chart configurations for the pruning rounds relative to every
ancestor.  \Cref{r:supported-products} replaces the interval-product
consequence exactly at its uses for missing direct, tensor,
Cartesian, and diagonal systems.  Thus each finite round either
loses only the bounded number of charts allowed by the matched
failure arguments or produces the next actual almost simple socle,
with all ancestor-relative exclusions.  These are the structural
conclusions of \Cref{c:classical-conclusion,c:alternating-conclusion}.

If a parent has bounded type, its bounded-order case is excluded by
residual height.  In bounded Lie rank, the bounded-equation pattern
argument of \Cref{sec:bounded} applies unchanged with a bounded
independent list: strict subgroup descents are strict after
intersection with the socle, whereas the spaces of equations have
bounded dimension.  Thus one obtains a sandwich throughout the
interval over a selected increment $t$.  Its height is $t$ and the
same divisor count gives at most $C2^t$ labels.  The exact
$t$-dimensional vector-subspace interval has at least
$q^{\lfloor t^2/4\rfloor}$ labels, a contradiction for sufficiently
large fixed $t$.  The sufficiently-general input is used here in
precisely the scope established in \Cref{sec:bounded}
\citep{LarsenPink2011}.

For a high-type path, use linear flags in the final protection
selection of \Cref{w:selection}.  All bad containing labels belong
to the exact subspace intervals just proved.  The same
exponential-gap estimate for flag proportions applies; hence the
closure-field tensor and root protections hold simultaneously for
the required ancestor comparisons.  The absolute group-theoretic
bound \Cref{w:path-bound} on the length of such a protected path
then contradicts a
choice of depth $D$ larger than that bound.

There is no dependence of the earlier finite choices on $n$.
First fix the high thresholds and the protected path bound, and
choose $D$ beyond it.  Next choose a prime $q$ above every absolute
line bound, then $r$ for the bounded-type count and protection
estimates.  Choose trial dimensions backwards through the finite
pruning rounds, allowing at each round all Hales--Jewett and
\Cref{r:affine-kernels} slot requirements for the trial and bounded
path ranks.  The latter requirements depend only on $q$, desired
fixed dimensions, and absolute orbit bounds.  The initial
comparisons require only a line or a bounded independent list, as
proved in \Cref{r:linear-patterns}; none requires comparison with
$|A_n|$.  Choose $N$ and the number of private charts large enough
for the finite color and chart-loss bounds.  Only then choose the
fields of \eqref{r:attachment-group} and finally $n$.
All elementary groups $X$ used are now bounded in size, so $n/|X|$
can exceed all normal-test and residual-height requirements while
$n$ accommodates every replicated regular action.  This gives one
fixed forbidden degree $n$.  For every $v\ge n$, the supported
$A_n$ restriction of \Cref{r:site-normal} gives the same contradiction.
\end{proof}
